% 第一部：作者、寫作脈絡、對話對象與本文的新意。
% 出處：刊本作者註（1346）、導論（1348–1353）；課程倉 notes/docs/1151/seminar/W06_1019_導讀問答.md 2026-10-06 兩則與其查核報告。
\chapter{作者、寫作脈絡與對話對象}

\qaentry{第一部之一}{作者 Jeremy Waldron}

\begin{qaquestion}
Waldron 撰文時任職何處？學術訓練與經歷為何？\\
這些背景在本文留下哪些痕跡？
\end{qaquestion}

\subsection{任職與經歷}

刊本的作者註寫明，Waldron 撰文時是哥倫比亞大學法學院的 University Professor，2006 年 7 月起轉任紐約大學法學院教授\pg{1346}。紐約大學法學院 2006 年的院刊介紹他時寫到：他在紐西蘭南部的 Invercargill 長大，在 Otago 大學取得哲學與法律學位，1978 年赴牛津攻讀研究所，1980 至 1982 年任牛津 Lincoln College 研究員，其後任教於愛丁堡大學、柏克萊、普林斯頓與哥倫比亞；同一篇介紹以 \en{taking strong, sometimes controversial stands on judicial review (should be weakened)} 概括他在違憲審查上的立場。\footcite{nyumag2006}

\subsection{紐西蘭背景在本文的痕跡}

紐西蘭之名在正文與註中出現於七頁（1349、1356、1361、1362、1365、1370、1394），用途分成兩組。

\begin{itemize}
\item 正面的例子。英國國會 1960 年代就墮胎、同性性行為、死刑所作的辯論，與加拿大、澳洲、紐西蘭等地的類似辯論，被用來反駁「議員無法負責任地處理這類問題」之說\pg{1349}；1990 年權利法案法是「更弱的形式」違憲審查的實例\pg{1356}；1993 年國會立法把選制改為比例代表制，說明立法機關能自行改革選舉制度\pg{1362~註~48}；權利法案法第 7 條要求檢察總長在法律案提出後指出其中與權利法案不符的條文，作者說這類安排是他所能接受的\pg{1370~註~63}。
\item 反面的例子。註 47 引他自己 2005 年在 New Zealand Law Journal 的短文，表示紐西蘭的一院制可能已使該國落在本文論證的範圍之外\pg{1361}；第柒部分註 136 再引這篇短文，作為論證不成立的一例\pg{1402}。
\end{itemize}

\begin{readernote}
作者用四項假設檢驗自己的國家，結論是紐西蘭未必合格。導讀時可以把這兩組並列：他所主張的制度在紐西蘭有實例，而他也承認紐西蘭的國會可能不符合第一項假設。
\end{readernote}

\qaentry{第一部之二}{作者為何在 2006 年再寫一篇反對違憲審查的文章}

\begin{qaquestion}
導論說雙方陣線已定，作者本人也寫過不少，為何再寫一篇？他自己給了哪些理由？
\end{qaquestion}

導論自問此事\pg{1350--1351}，答以四項理由。

\begin{enumerate}
\item 找出獨立於歷史表現與個案結果的核心論證。他寫到，Tushnet 與 Kramer 的近著把對違憲審查的理論批評，與它的歷史起源、以及兩人對一部較少司法化的美國憲法的構想交纏在一起，並表明這不是對兩人的批評。Michelman 為 Kramer 一書所寫的推薦語說，該書的歷史 \en{puts flesh on the bones of debates over judicial review}；Waldron 接著說，他要 \en{take off some of the flesh and boil down the normative argument to its bare bones}\pg{1351}。
\item 讓論辯脫離對個別判決的好惡。他引 Charles Black 的話，人們對違憲審查的反對往往是 \en{a sometime thing}，並舉墮胎權論辯為例：支持選擇權的一方擔心動搖 Roe v.\ Wade，連反對違憲審查的論證都不願考慮\pg{1351--1352}。
\item 讓論證適用於美國以外的國家。英國人討論本國法官有限的審查權時，不關心 1805 年聯邦黨人與共和黨人的爭論，也不關心 Brown 的遺產\pg{1352}。
\item 對自己先前的表述不滿意。他寫道，自己過去的討論把違憲審查與法理學、政治哲學的其他問題交纏在一起，未曾清楚陳述基本的反對理由，也未充分回應對 Law and Disagreement 等著作的批評\pg{1352--1353}。
\end{enumerate}

摘要把全文的主張歸成兩項：其一，沒有理由認為權利在違憲審查之下比在民主立法機關之下受到更好的保護；其二，撇開所產生的結果，違憲審查在民主上欠缺正當性。摘要自述第二項為人熟知，第一項較少人提出\pg{1346}。導論末段把同一主張寫成兩條戰線：違憲審查以先例、文本與解釋的枝節問題分散公民對真正爭點的注意；它讓少數非民選、不須負責的法官以多數決作終局決定，剝奪一般公民的參政權，撇開代表與政治平等的原則\pg{1353}。

\qaentry{第一部之三}{寫作時的論辯脈絡}

\begin{qaquestion}
本文發表時，美國與其他普通法國家的違憲審查論辯處於什麼狀態？本文在三次發表場合面對的是哪些人？
\begin{enumerate}
\item 美國的論辯。
\item 大英國協的新模式。
\item 發表場合與評論人。
\end{enumerate}
\end{qaquestion}

\begin{qafigure}{本文前後的相關事件。年份與刊本頁的出處見本則各小節。\label{fig:timeline}}
% 本文前後的相關事件：蛇行時間軸，每列四則（改自 ~/.claude/skills/diagram-design/templates/蛇行時間軸.tex）。
% 灰階為底，本文一則用強調色。事件卡第一行小字為類別。
\begin{tikzpicture}[qa,
    track/.style={line width=1pt,color=ruleink,line cap=round},
    dot/.style={circle,fill=#1,draw=white,line width=1pt,inner sep=0pt,minimum size=6.5pt},
    dt/.style={font=\footnotesize\bfseries,text=#1,inner sep=1pt},
    card/.style={nbox,text width=29mm,font=\scriptsize,text=caseink,anchor=north,
      execute at begin node={\raggedright\setstretch{1.2}}},
    keycard/.style={card,kbox}]
  \def\xa{1.9}\def\dx{3.4}\def\dy{3.0}\def\R{1.5}
  \pgfmathsetmacro{\xl}{\xa-0.7}\pgfmathsetmacro{\xr}{\xa+3*\dx+0.7}
  \draw[track] (\xl,0) -- (\xr,0) arc[start angle=90,end angle=-90,radius=\R]
    -- (\xl,-\dy) arc[start angle=90,end angle=270,radius=\R]
    -- (\xr,-2*\dy) arc[start angle=90,end angle=-90,radius=\R]
    -- ({\xa+1.6*\dx},-3*\dy);
  \draw[track,-{Stealth[length=2.4mm,width=2.2mm]}] ({\xa+1.6*\dx},-3*\dy) -- ({\xa+1.4*\dx},-3*\dy);
  \foreach \r/\c/\k/\d/\t in {
      0/0/n/1938/{\textcolor{dimink}{判決}\\Carolene Products 案腳註四：孤立而隔絕的少數},
      0/1/n/1962/{\textcolor{dimink}{著作}\\Bickel, The Least Dangerous Branch：對抗多數的難題},
      0/2/n/1982/{\textcolor{dimink}{立法}\\加拿大權利與自由憲章第 33 條},
      0/3/n/1990/{\textcolor{dimink}{立法}\\紐西蘭權利法案法},
      1/0/n/1998/{\textcolor{dimink}{立法}\\英國人權法：不相容宣告},
      1/1/n/1999/{\textcolor{dimink}{著作}\\Waldron, Law and Disagreement；Tushnet, Taking the Constitution Away from the Courts},
      1/2/n/2003/{\textcolor{dimink}{判決}\\Goodridge 案：麻州同性婚姻},
      1/3/n/2004/{\textcolor{dimink}{著作}\\Kramer, The People Themselves},
      2/0/k/2006/{\textcolor{dimink}{本文}\\115 Yale L.J. 1346},
      2/1/n/2008/{\textcolor{dimink}{著作}\\Fallon 以回應本文為題},
      2/2/n/2014/{\textcolor{dimink}{著作}\\Waldron, Five to Four：法院內部的多數決},
      2/3/n/2015/{\textcolor{dimink}{判決}\\Obergefell 案，五比四},
      3/0/n/2022/{\textcolor{dimink}{判決}\\Dobbs 案：推翻 Roe 與 Casey},
      3/1/n/2022/{\textcolor{dimink}{立法}\\紐西蘭權利法案法增訂第 7A 條}}{
    \pgfmathtruncatemacro{\odd}{mod(\r,2)}
    \pgfmathsetmacro{\x}{\odd==0 ? \xa+\c*\dx : \xa+(3-\c)*\dx}
    \pgfmathsetmacro{\y}{-\r*\dy}
    \if k\k
      \node[dot=figE] at (\x,\y) {};
      \node[dt=figE,anchor=south] at (\x,\y+0.16) {\d};
      \node[keycard] at (\x,\y-0.22) {\t};
    \else
      \node[dot=caseink] at (\x,\y) {};
      \node[dt=caseink,anchor=south] at (\x,\y+0.16) {\d};
      \node[card] at (\x,\y-0.22) {\t};
    \fi
  }
\end{tikzpicture}

\end{qafigure}

\subsection{美國的論辯}

導論描述了 2006 年前後的美國論辯\pg{1350}。長期以來，支持違憲審查的是自由派，反對的是不滿自由派法院所支持之權利的保守派；Rehnquist 法院廢棄若干自由派立法之後（註~12 舉 United States v.\ Morrison、United States v.\ Lopez），自由派的反對也增加了。註 11 列 Kramer 的 The People Themselves（2004）與 Tushnet 的 Taking the Constitution Away from the Courts（1999）為近年攻擊違憲審查之書，Dworkin、Eisgruber、Sager 則為違憲審查辯護（註~\mbox{11、13}）。Marbury v.\ Madison 兩百週年與 Brown v.\ Board of Education 五十週年，又帶出一批關於違憲審查起源與成就的討論。

\subsection{大英國協的新模式}

第壹部分界定強式與弱式違憲審查時，除美國與部分歐洲法院之外，另舉大英國協的三例：英國 1998 年人權法的不相容宣告為弱式\pg{1355}，紐西蘭 1990 年權利法案法只要求法院儘可能作合於權利的解釋，為更弱的形式\pg{1356}，加拿大權利與自由憲章第 33 條的「儘管條款」雖使立法機關得排除審查，作者仍計為強式\pg{1356--1357}。本文的攻擊對象限於強式審查，前兩種制度提供了作者所能接受的對照（見第三部）；加拿大屬於攻擊的標的。

\subsection{發表場合與評論人}

作者註列出本文的三次發表：倫敦大學學院的法律與社會哲學研討會、耶路撒冷希伯來大學的法學院工作坊、哈佛法學院的憲法研討會，Dworkin、Ruth Gavison、Seana Shiffrin 擔任正式評論人；致謝名單中還有 Raz、Fallon、Michelman 等人\pg{1346}。正文所回應的論者，有幾位就在這份名單上。

\qaentry{第一部之四}{本文與作者先前著作的關係}

\begin{qaquestion}
本文是 Waldron 違憲審查理論的首次陳述，還是重述？與 Law and Disagreement 等書是什麼關係？
\end{qaquestion}

註 14 列出他此前討論違憲審查的著作，包括 Law and Disagreement（1999）第 10 至 17 頁與第 211 至 312 頁、Deliberation, Disagreement, and Voting（1999）、Judicial Power and Popular Sovereignty（2002）；註 19 另列他從權利理論、公民共和主義、盧梭與邊沁兩種民主觀、歐陸人民主權論等角度處理同一問題的文章\pg{1351--1353}。\footcite{waldron1999lad} 第肆部分註 89 引 The Dignity of Legislation（1999），說明他刻意以樂觀的立法機關圖像對應法院支持者的樂觀圖像\pg{1379}。\footcite{waldron1999dignity}

依作者自述\pg{1352--1353}，本文重述他已有的立場，把論證從上列著作與政治哲學的其他問題中抽出，寫成可以單獨檢驗的形式。四項假設、結果面與程序面兩類理由的區分、核心情形與非核心情形的劃分，都在本文寫成可以單獨檢驗的形式；註 42 說明四項假設改寫自作者 2004 年 Some Models of Dialogue Between Judges and Legislators 一文。

\qaentry{第一部之五}{本文和誰對話}

\begin{qaquestion}
本文引用、回應了哪些論者？依立場可分幾群？作者對各說如何回應？
\end{qaquestion}

\begin{qafigure}{本文的對話對象，依作者對各人的態度分三組。長條為本文正文與註中提及該姓氏的次數，取自課程倉 \texttt{literature-mentions.tsv}。\label{fig:interlocutors}}
% 本文的對話對象：三組，依位置分組，灰階；長條長度＝本文提及次數（literature-mentions.tsv）。
\begin{tikzpicture}[qa]
  \def\u{0.12}% 每次提及 1.2mm
  \newcommand{\who}[4]{% x, y, 名, 次數（0 表示未計）
    \node[anchor=east,font=\footnotesize] at (#1,#2) {#3};
    \ifnum#4>0
      \fill[figGray] (#1+0.1,#2-0.07) rectangle ++({#4*\u},0.14);
      \node[lbl,anchor=west] at ({#1+0.15+#4*\u},#2) {#4};
    \fi}
  \node[hdr,anchor=west] at (0,0.6) {所反駁的辯護者};
  \who{2.2}{0}{Dworkin}{28}
  \who{2.2}{-0.45}{Rawls}{17}
  \who{2.2}{-0.9}{Raz}{11}
  \who{2.2}{-1.35}{Eisgruber}{7}
  \who{2.2}{-1.8}{Moore}{0}
  \who{2.2}{-2.25}{Christiano}{0}
  \node[hdr,anchor=west] at (8.4,0.6) {同樣批評違憲審查};
  \who{10.6}{0}{Tushnet}{10}
  \who{10.6}{-0.45}{Kramer}{8}
  \node[hdr,anchor=west] at (8.4,-0.95) {借用其概念或論證工具};
  \who{10.6}{-1.4}{Michelman}{6}
  \who{10.6}{-1.85}{Bickel}{3}
  \node[lbl,anchor=north west,text width=60mm,align=left] at (8.4,-2.15) {Kavanagh、Fuller、Wollheim、Condorcet、Elster、Scalia（Casey 案不同意見）};
\end{tikzpicture}

\end{qafigure}

\subsection{作者所反駁的辯護者}

\begin{longtable}{@{}L{0.17\linewidth}L{0.36\linewidth}L{0.39\linewidth}@{}}
\toprule
論者 & 主張 & 作者的回應 \\
\midrule
\endhead
Dworkin & 法院是「原則的論壇」\pg{1385--1386}；參與的價值在於確認公民平等成員地位的後果\pg{1375}；尊重權利是任何決策制度具正當性的背景條件\pg{1398--1400} & 法院的推理以解釋與教義為主；Dworkin 嚴重低估了參與權的觀念；「彼此尊重權利」有兩種讀法，作者採「認真對待並真誠設法確認」一讀，在核心情形下已由第三項假設滿足 \\
Raz & 結果面理由是唯一值得考慮的理由\pg{1374}；立法機關受宗派利益驅使\pg{1376--1378} & 前者被稱為教條主義；同樣的宗派壓力也見於法院（Korematsu 等案），立法機關若普遍如此，即與第三項假設不相容，屬非核心情形 \\
Moore & 法官每天面對具體個人的處境，較能得到道德洞見\pg{1379} & 案件上訴到最高審級時，原告的個人處境已幾乎不見，爭的是抽象的權利問題\pg{1380} \\
Eisgruber & 法官經提名同意而有民主資格\pg{1394~註~125} & 法官的民主資格遠不及議員 \\
Rawls & 良序社會亦宜設立違憲審查\pg{1366~註~53} & 作者的目標之一是證明 Rawls 對良序社會與違憲審查的看法有誤；同時借用 Rawls 的「判斷的負擔」作為第四項假設的依據\pg{1368} \\
\bottomrule
\end{longtable}

\subsection{同樣批評違憲審查的近著}

Tushnet 與 Kramer 被列為近年攻擊違憲審查的著作\pg{1350~註~11}。作者表明不批評兩書，只把其中的歷史與制度構想抽去，留下規範論證\pg{1351}。

\subsection{借用其概念或論證工具者}

\begin{itemize}
\item Bickel 的「對抗多數的難題」與「脫軌制度」，作為既有反對論的代表\pg{1349}。\footcite[16--18]{bickel1962}
\item Michelman 所提的問法（何種方法最可能觸及權利真相、同時充分尊重受影響者平等發言的主張），作者將其前後倒置成第二種問法，用以設定比較兩類理由的問題\pg{1375~註~76}。
\item Kavanagh：結果面理由不必預設權利內容，只問何種程序較可能觸及權利的真相\pg{1373--1374}。
\item Fuller 與 Hart、Sacks 的法律程序學派，作為第肆部分檢驗機構能力的起點\pg{1376}。
\item Wollheim 的民主理論弔詭，作者改說成一切政治理論都有的弔詭，用以說明任何決定程序都會出錯\pg{1372}；Christiano 的程序無窮後退則是作者所回應的反對意見，作者指 Christiano 未證明其為惡性後退\pg{1371~註~66}。
\item Condorcet 陪審團定理，作者指出它反而有利於人數較多的立法機關\pg{1392}。
\item Elster 的 Ulysses 之喻，見於五項調和論證中的預先承諾說\pg{1393--1394}。
\item Scalia 在 Casey 案的不同意見，作者據以說明法院的正當性問題首當其衝\pg{1390--1391}。
\end{itemize}

\qaentry{第一部之六}{本文的新意}

\begin{qaquestion}
與 Bickel 以來的反對論相比，本文提出了哪些新的論證工具？各在何處？
\end{qaquestion}

摘要自述，主張違憲審查不能更好地保護權利的那一項較少人提出\pg{1346}。依作者在本文所用的工具，可整理成下表。

\begin{longtable}{@{}L{0.24\linewidth}L{0.58\linewidth}L{0.10\linewidth}@{}}
\toprule
工具 & 內容 & 刊本頁 \\
\midrule
\endhead
核心論據的寫法 & 抽離歷史表現與個案結果，只留規範論證 & 1351 \\
附條件的論證 & 論證以四項假設為前提；假設全部成立為核心情形，任一不成立屬非核心情形，論證可能不成立 & 1359--1360 \\
以權利分歧為出發點 & 社會成員認真對待權利，而就有哪些權利、其內容為何存在善意的實質分歧；與道德相對主義無關 & 1366--1369 \\
兩類理由 & 結果面問何者較可能得出正確結果；程序面問何人有權參與決定，與結果無涉 & 1372--1373 \\
C\textsubscript{n} 的兩個問題 & 向不同意決定的公民說明：為何由這些人決定、為何採此決定程序 & 1387 \\
法院內部的多數決 & 大法官之間的簡單多數決欠缺公平性的論證 & 1391--1392 \\
兩種多數與少數 & 決定上的少數與議題上的少數重合時，始構成多數暴政 & 1397--1398 \\
\bottomrule
\end{longtable}

\begin{readernote}
前三項決定了論證的形狀：先限定爭點，再限定社會，最後只比較兩個機關。後四項是比較時用的尺。導讀時可以先講前三項，聽者才知道後面的比較在什麼條件下成立。
\end{readernote}
